\section*{Resumo}
Este trabalho apresenta o projeto e a validação de um rastreador de referência para um modelo farmacocinético discreto de três compartimentos, considerando incertezas paramétricas na dinâmica da planta. O sistema utiliza um controlador proporcional-integral (PI), cujos ganhos foram sintetizados por meio de um procedimento iterativo baseado em desigualdades matriciais lineares (LMIs). Para representar as incertezas, adotou-se um politopo com oito vértices, definidos a partir dos valores extremos de três parâmetros do modelo. Inicialmente, simulou-se o comportamento do sistema em malha fechada no Matlab/Simulink para os oito vértices, utilizando os ganhos previamente sintetizados do controlador PI. Na sequência, desenvolveu-se uma plataforma de emulação embarcada composta por dois microcontroladores ESP32, um representando o controlador e o outro, a planta. Os dispositivos se comunicaram por UART, enquanto o protocolo MQTT e uma interface Node-RED permitiram modificar a referência de entrada e acompanhar o estado do sistema em tempo real. Tanto na simulação quanto na emulação, a saída da planta se aproximou da referência em todos os vértices avaliados, embora tenham sido observadas diferenças no comportamento transitório durante os ensaios. Por meio de potenciômetros conectados ao ESP32 responsável pela planta, realizaram-se testes variando os parâmetros incertos dentro da região do politopo. Esses testes mostraram que a malha permaneceu estável nas condições avaliadas, sem necessidade de reajustar os ganhos do controlador. Dessa forma, os resultados demonstram a viabilidade da plataforma desenvolvida, dado que o controlador manteve o rastreamento da referência e a estabilidade nas condições experimentais avaliadas.

\noindent\textbf{Palavras-chave:} controle robusto; desigualdades matriciais lineares; farmacocinética; rastreamento de referência; sistemas embarcados.
\newpage